% GENERATED FILE. Edit canonical lessons, then run npm run generate:books.
% canonical-source-hash: fnv1a64:3115728bdcd0a5fd
% canonical-lessons: RU-C179-retsept, RU-C179-ukol, RU-C179-rana, RU-C179-sinyak, RU-C179-perelom

\chapter{Prescription, Injection, Wound, Bruise, Fracture}
\label{ch:ru-prescription-injection-wound-bruise-fracture}
\hlchaptermodality{179}

\begin{chapteropening}
\textbf{By the end of this chapter:} \emph{I can say a prescription, an injection, a wound, a bruise and a fracture.}

Complete the last lesson of chapter 179: I can say a prescription, an injection, a wound, a bruise and a fracture.
\end{chapteropening}

\section[\texorpdfstring{\ru{рецепт} (retsépt)}{рецепт (retsépt)}]{\ru{рецепт} (retsépt) — a prescription; a recipe}

\label{lesson:RU-C179-retsept}



\begin{quote}
Before the new one: say the Russian for to go for a ride, then the Russian for to sunbathe.
\end{quote}

\subsection*{You'll want to know: \ru{рецепт}}
\textbf{\ru{рецепт}} (\emph{retsépt}) — \textquotedblleft{}a prescription; a recipe\textquotedblright{}.

\begin{grammarlens}[title={What you've built}]
One more word for this chapter.
\end{grammarlens}

\subsection*{Your turn}
\begin{itemize}
\raggedright
  \item \emph{Say it:} \emph{retsépt}
  \item \emph{Say it:} \emph{retsépt}, once more
  \item \emph{Say it:} say \emph{retsépt}
  \item \emph{Recall:} say the Russian for to go for a ride, then the Russian for to sunbathe, then say \emph{retsépt} again
\end{itemize}

\begin{tcolorbox}[breakable,colback=teal!4,colframe=teal!35!black,title={Before you move on}]
What does \ru{рецепт} mean? (\textquotedblleft{}A prescription; a recipe\textquotedblright{}.) Say it once more.
\end{tcolorbox}
\section[\texorpdfstring{\ru{укол} (ukól)}{укол (ukól)}]{\ru{укол} (ukól) — an injection}

\label{lesson:RU-C179-ukol}



\begin{quote}
Before the new one: say the Russian for to sunbathe, then the Russian for a prescription; a recipe.
\end{quote}

\subsection*{You'll want to know: \ru{укол}}
\textbf{\ru{укол}} (\emph{ukól}) — \textquotedblleft{}an injection\textquotedblright{}.

\begin{grammarlens}[title={What you've built}]
One more word for this chapter.
\end{grammarlens}

\subsection*{Your turn}
\begin{itemize}
\raggedright
  \item \emph{Say it:} \emph{ukól}
  \item \emph{Say it:} \emph{ukól}, once more
  \item \emph{Say it:} say \emph{ukól}
  \item \emph{Recall:} say the Russian for to sunbathe, then the Russian for a prescription; a recipe, then say \emph{ukól} again
\end{itemize}

\begin{tcolorbox}[breakable,colback=teal!4,colframe=teal!35!black,title={Before you move on}]
What does \ru{укол} mean? (\textquotedblleft{}An injection\textquotedblright{}.) Say it once more.
\end{tcolorbox}
\section[\texorpdfstring{\ru{рана} (rána)}{рана (rána)}]{\ru{рана} (rána) — a wound}

\label{lesson:RU-C179-rana}



\begin{quote}
Before the new one: say the Russian for a prescription; a recipe, then the Russian for an injection.
\end{quote}

\subsection*{You'll want to know: \ru{рана}}
\textbf{\ru{рана}} (\emph{rána}) — \textquotedblleft{}a wound\textquotedblright{}.

\begin{grammarlens}[title={What you've built}]
One more word for this chapter.
\end{grammarlens}

\subsection*{Your turn}
\begin{itemize}
\raggedright
  \item \emph{Say it:} \emph{rána}
  \item \emph{Say it:} \emph{rána}, once more
  \item \emph{Say it:} say \emph{rána}
  \item \emph{Recall:} say the Russian for a prescription; a recipe, then the Russian for an injection, then say \emph{rána} again
\end{itemize}

\begin{tcolorbox}[breakable,colback=teal!4,colframe=teal!35!black,title={Before you move on}]
What does \ru{рана} mean? (\textquotedblleft{}A wound\textquotedblright{}.) Say it once more.
\end{tcolorbox}
\section[\texorpdfstring{\ru{синяк} (sinyák)}{синяк (sinyák)}]{\ru{синяк} (sinyák) — a bruise}

\label{lesson:RU-C179-sinyak}



\begin{quote}
Before the new one: say the Russian for an injection, then the Russian for a wound.
\end{quote}

\subsection*{You'll want to know: \ru{синяк}}
\textbf{\ru{синяк}} (\emph{sinyák}) — \textquotedblleft{}a bruise\textquotedblright{}.

\begin{grammarlens}[title={What you've built}]
One more word for this chapter.
\end{grammarlens}

\subsection*{Your turn}
\begin{itemize}
\raggedright
  \item \emph{Say it:} \emph{sinyák}
  \item \emph{Say it:} \emph{sinyák}, once more
  \item \emph{Say it:} say \emph{sinyák}
  \item \emph{Recall:} say the Russian for an injection, then the Russian for a wound, then say \emph{sinyák} again
\end{itemize}

\begin{tcolorbox}[breakable,colback=teal!4,colframe=teal!35!black,title={Before you move on}]
What does \ru{синяк} mean? (\textquotedblleft{}A bruise\textquotedblright{}.) Say it once more.
\end{tcolorbox}
\section[\texorpdfstring{\ru{перелом} (perelóm)}{перелом (perelóm)}]{\ru{перелом} (perelóm) — a fracture}

\label{lesson:RU-C179-perelom}



\begin{quote}
Before the new one: say the Russian for a wound, then the Russian for a bruise.
\end{quote}

\subsection*{You'll want to know: \ru{перелом}}
\textbf{\ru{перелом}} (\emph{perelóm}) — \textquotedblleft{}a fracture\textquotedblright{}.

\begin{grammarlens}[title={What you've built}]
One more word for this chapter.
\end{grammarlens}

\subsection*{Your turn}
\begin{itemize}
\raggedright
  \item \emph{Say it:} \emph{perelóm}
  \item \emph{Say it:} \emph{perelóm}, once more
  \item \emph{Say it:} say \emph{perelóm}
  \item \emph{Recall:} say the Russian for a wound, then the Russian for a bruise, then say \emph{perelóm} again
\end{itemize}

\begin{tcolorbox}[breakable,colback=teal!4,colframe=teal!35!black,title={Before you move on}]
What does \ru{перелом} mean? (\textquotedblleft{}A fracture\textquotedblright{}.) Say it once more.
\end{tcolorbox}
